\Askhsh{36840_SOLUTION}{1}
\begin{ylh}{1}
    \item [Ύλη από εικόνα]
\end{ylh}
\begin{center}
    \textbf{\Large ΛΥΣΗ}
\end{center}

\begin{alist}
    \item Το σύνολο των τετμημένων των σημείων της $C_f$ και της $C_h$ αντίστοιχα, αποτελεί το πεδίο ορισμού της κάθε συνάρτησης. Από τις γραφικές παραστάσεις του σχήματος παρατηρούμε ότι το πεδίο ορισμού της συνάρτησης $f$ είναι το διάστημα $[2,7]$ και το πεδίο ορισμού της συνάρτησης $h$ είναι το διάστημα $[5,7]$.
    \item Παρατηρούμε ότι και οι 2 γραφικές παραστάσεις των $f$ και $h$ έχουν κοινό σημείο με τον άξονα $x'x$ το σημείο $A(6,0)$, οπότε ισχύει $f(6) = h(6) = 0$.
    \begin{rlist}
        \item $\lim_{x \to 6} \frac{1}{f(x)} = +\infty$ γιατί $\lim_{x \to 6} f(x) = f(6) = 0$ και $f(x) > 0$ κοντά στο 6, αφού η $C_f$ εφάπτεται στον άξονα $x'x$ και βρίσκεται πάνω από αυτόν κοντά στο 6.
        \item Το $\lim_{x \to 6} \frac{1}{h(x)}$ δεν υπάρχει. Γιατί:
        \[ \lim_{x \to 6} h(x) = h(6) = 0, \text{ αφού η } C_h \text{ εφάπτεται στον άξονα } x'x \text{ στο σημείο } A(6,0). \]
        Είναι $h(x) < 0$ για $x < 6$. Άρα $\lim_{x \to 6^-} \frac{1}{h(x)} = -\infty$.
        Ενώ $h(x) > 0$ για $x > 6$. Άρα $\lim_{x \to 6^+} \frac{1}{h(x)} = +\infty$.
        \item Για το $\lim_{x \to 6} \frac{f(x)}{x-6}$ έχουμε:
        \[ \lim_{x \to 6} \frac{f(x)}{x-6} = \lim_{x \to 6} \frac{f(x)-0}{x-6} = \lim_{x \to 6} \frac{f(x)-f(6)}{x-6} = f'(6), \]
        αφού η $f$ από την υπόθεση είναι παραγωγίσιμη συνάρτηση στο πεδίο ορισμού της, άρα και στο $x_0 = 6$.
        Όμως, από την υπόθεση, η γραφική παράσταση της $f$ δέχεται οριζόντια εφαπτομένη στο σημείο της $A(6,0)$, τον άξονα $x'x$. Οπότε, η παράγωγος της στο σημείο αυτό, δηλαδή το $f'(6)$, θα ισούται με 0.
        Άρα, τελικά $\lim_{x \to 6} \frac{f(x)}{x-6} = 0$.
    \end{rlist}
\end{alist}